\documentclass[12pt]{article}
\usepackage[T1]{fontenc}
\usepackage[utf8]{inputenc}
\usepackage{lmodern}
\usepackage{textcomp}
\usepackage{enumitem}
\usepackage{geometry}
\geometry{margin=1in}
\begin{document}
\begin{center}
Answer Key - Computer Science - Graduate Level Assessment 12 \
\small (Answers are ordered exactly as in the question paper.)
\end{center}

\section*{Multiple Choice}
\begin{enumerate}[label=\arabic*.]
\item \textbf{Answer:} C
\\ \textit{Options:} A) They are generated when memory cycles are "stolen".; B) They are used in place of data channels.; C) They can indicate completion of an I/O operation.; D) They cannot be generated by arithmetic operations.
\\ \textit{Note:} Correct option: C - They can indicate completion of an I/O operation.
\item \textbf{Answer:} B
\\ \textit{Options:} A) Installing and configuring a IDS that can read the IP header; B) Comparing the TTL values of the actual and spoofed addresses; C) Implementing a firewall to the network; D) Identify all TCP sessions that are initiated but does not complete successfully
\\ \textit{Note:} Correct option: B - Comparing the TTL values of the actual and spoofed addresses
\item \textbf{Answer:} B
\\ \textit{Options:} A) No, I cannot compute the key.; B) Yes, the key is k = m xor c.; C) I can only compute half the bits of the key.; D) Yes, the key is k = m xor m.
\\ \textit{Note:} Correct option: B - Yes, the key is k = m xor c.
\item \textbf{Answer:} C
\\ \textit{Options:} A) Replaced; B) Overview; C) Changed; D) Violated
\\ \textit{Note:} Correct option: C - Changed
\item \textbf{Answer:} C
\\ \textit{Options:} A) Active, inactive, standby; B) Open, half-open, closed; C) Open, filtered, unfiltered; D) Active, closed, unused
\\ \textit{Note:} Correct option: C - Open, filtered, unfiltered
\end{enumerate}

\section*{True / False}
\begin{enumerate}[label=\arabic*.]
\setcounter{enumi}{5}
\item \textbf{Answer:} False
\\ \textit{Options:} A) True; B) False
\\ \textit{Note:} LLM-generated false statement. Correct answer: 'Snort'.
\item \textbf{Answer:} False
\\ \textit{Options:} A) True; B) False
\\ \textit{Note:} LLM-generated false statement. Correct answer: 'CBF'.
\item \textbf{Answer:} False
\\ \textit{Options:} A) True; B) False
\\ \textit{Note:} LLM-generated false statement. Correct answer: 'EtterPeak'.
\item \textbf{Answer:} False
\\ \textit{Options:} A) True; B) False
\\ \textit{Note:} LLM-generated false statement. Correct answer: 'Freenet'.
\item \textbf{Answer:} True
\\ \textit{Options:} A) True; B) False
\\ \textit{Note:} LLM-generated true statement based on MCQ.
\end{enumerate}

\section*{Long Form Response}
\begin{enumerate}[label=\arabic*.]
\setcounter{enumi}{10}
\item \textbf{Answer:} def find\_literals(text, pattern): | return (match.re.pattern, s, e)
\\ \textit{Note:} Students should provide a detailed explanation referencing algorithm steps.
\item \textbf{Answer:} def tn\_gp(a,n,r): | return tn
\\ \textit{Note:} Students should provide a detailed explanation referencing algorithm steps.
\end{enumerate}

\vfill
\noindent\textit{End of Answer Key}
\end{document}